\documentclass[11pt,a4paper]{article}

% ==========================================================
% Configuración general
% ==========================================================
\usepackage[utf8]{inputenc}
\usepackage[T1]{fontenc}
\usepackage{lmodern}
\usepackage{microtype}
\usepackage[spanish]{babel}
\usepackage[margin=2.1cm]{geometry}
\usepackage{graphicx}
\usepackage{booktabs}
\usepackage{array}
\usepackage{tabularx}
\usepackage{longtable}
\usepackage{amssymb}
\usepackage{tikz}
\usetikzlibrary{arrows.meta,positioning}
\usepackage{float}
\usepackage{hyperref}
\usepackage{xurl}
\usepackage{xcolor}
\usepackage{enumitem}
\usepackage{titlesec}
\usepackage{fancyhdr}
\emergencystretch=2em

% Paleta de color (misma que la presentación en beamer)
\definecolor{myNewColorA}{RGB}{27,94,55}
\definecolor{myNewColorB}{RGB}{46,125,80}
\definecolor{myNewColorC}{RGB}{184,218,191}

\titleformat{\section}{\Large\bfseries\color{myNewColorA}\raggedright}{\thesection.}{0.5em}{}
\titleformat{\subsection}{\large\bfseries\color{myNewColorB}\raggedright}{\thesubsection}{0.5em}{}

\hypersetup{
    colorlinks=true,
    linkcolor=myNewColorA,
    urlcolor=myNewColorA,
    citecolor=myNewColorA,
    pdftitle={Avance 1 - Proyecto Integrador FireForest},
    pdfauthor={Malan Mullo, Marco Vinicio; Pujos Culque, Viviana Isabel; Chuncho Morocho, Carlos Guillermo}
}

\pagestyle{fancy}
\fancyhf{}
\renewcommand{\headrulewidth}{0.4pt}
\fancyhead[L]{\small Proyecto Integrador -- Avance 1}
\fancyhead[R]{\small FireForest -- Grupo 4}
\fancyfoot[C]{\small\thepage}

\setlist[itemize]{itemsep=1pt, topsep=1pt}
\setlist[enumerate]{itemsep=1pt, topsep=1pt}

% Columnas de párrafo con salto de línea controlado (para tablas)
\newcolumntype{L}[1]{>{\raggedright\arraybackslash}p{#1}}
\newcolumntype{C}[1]{>{\centering\arraybackslash}p{#1}}

\begin{document}

% ==========================================================
% PORTADA
% ==========================================================
\begin{titlepage}
\centering
\vspace*{1cm}

{\large ESCUELA SUPERIOR POLIT\'ECNICA DE CHIMBORAZO (ESPOCH)\par}
\vspace{0.3cm}
{\large Maestr\'ia en Estad\'istica con menci\'on en Ciencia de Datos e Inteligencia Artificial\par}

\vspace{2.5cm}
{\Huge\bfseries\color{myNewColorA} Avance del Proyecto Integrador\par}
\vspace{0.5cm}
{\large Problema, fuentes, arquitectura preliminar y plan de implementaci\'on\par}

\vspace{1.5cm}
\begin{minipage}{0.85\textwidth}
\centering
{\Large\bfseries T\'itulo del proyecto\par}
\vspace{0.3cm}
{\large Dise\~no de una arquitectura h\'ibrida SQL--NoSQL para la integraci\'on y consulta de datos de incendios forestales en el cant\'on Loja\par}
\end{minipage}

\vspace{2cm}
{\large\bfseries Integrantes\par}
\vspace{0.3cm}
{\large
Malan Mullo, Marco Vinicio \\
Pujos Culque, Viviana Isabel \\
Chuncho Morocho, Carlos Guillermo \par}

\vspace{2cm}
{\large\bfseries Grupo 4}\par

\vspace{1.5cm}
{\large Avance 1 \par}
{\large 2026 \par}

\vfill
\end{titlepage}

% ==========================================================
% 1. DEFINICIÓN Y JUSTIFICACIÓN DEL PROBLEMA
% ==========================================================
\section{Definici\'on y justificaci\'on del problema}

Los incendios forestales en el cant\'on Loja se concentran en temporadas secas y afectan ecosistemas, biodiversidad y comunidades rurales cercanas a zonas de bosque y pastizal. El cant\'on se ubica en los Andes tropicales del sur de Ecuador, donde el relieve monta\~noso genera transiciones ambientales pronunciadas: temperatura, precipitaci\'on, humedad, vegetaci\'on y combustibles cambian con la altitud. Balaguer-Beser et al. (2026) documentan esta heterogeneidad en tres zonas de estudio del sur de Ecuador (Malacatos, Punzara y San Lucas, entre aproximadamente 1549 y 2757 m s.\,n.\,m.), con condiciones contrastantes de comportamiento del fuego entre ellas; estos resultados corresponden espec\'ificamente a esas tres zonas y se utilizan aqu\'i para contextualizar la heterogeneidad ambiental y altitudinal del sur del Ecuador, no como una descripci\'on generalizable a la totalidad del cant\'on Loja.

La informaci\'on necesaria para estudiar estos eventos proviene de fuentes muy distintas entre s\'i: detecciones satelitales de focos de calor (VIIRS), registros clim\'aticos de precipitaci\'on (CHIRPS) y una malla espacial de referencia del territorio. Cada fuente tiene su propio formato, resoluci\'on y frecuencia de actualizaci\'on, y hoy no existe una forma sencilla de integrarlas para responder preguntas conjuntas. Este es, en esencia, un problema de \textbf{Ingenier\'ia de Datos}: dise\~nar el recorrido completo de los datos --- captura, almacenamiento, integraci\'on y consolidaci\'on --- antes de poder analizarlos.

El sistema involucrado es el de monitoreo y an\'alisis de riesgo de incendios forestales a nivel territorial, en el que intervienen t\'ipicamente actores como los Gobiernos Aut\'onomos Descentralizados (GAD Municipal y GAD Provincial), el Cuerpo de Bomberos y el Ministerio del Ambiente, Agua y Transici\'on Ecol\'ogica (MAATE), as\'i como las comunidades rurales asentadas cerca de zonas boscosas del cant\'on. FireForest es un ejercicio acad\'emico: no existe un convenio con estas instituciones, se mencionan como referencia del tipo de actor que se beneficiar\'ia de una soluci\'on de este tipo, y el proyecto trabaja \'unicamente con datos de ejemplo, muestras controladas o simulados. Los usuarios o beneficiarios potenciales son gestores de riesgo y personal t\'ecnico de instituciones ambientales o de emergencia, investigadores y estudiantes que analicen patrones espacio-temporales de incendios forestales, y comunidades rurales interesadas en conocer zonas de mayor recurrencia cerca de sus localidades.

\textbf{Objetivo general:} dise\~nar una arquitectura h\'ibrida SQL--NoSQL para integrar datos de detecciones VIIRS, precipitaci\'on CHIRPS y una malla espacial del cant\'on Loja, con el prop\'osito de construir un dataset anal\'itico reproducible por celda y mes.

\textbf{Pregunta principal:} \textquestiondown c\'omo dise\~nar una arquitectura h\'ibrida SQL--NoSQL que permita integrar, almacenar y consultar de forma reproducible datos heterog\'eneos de incendios forestales en el cant\'on Loja?

\textbf{Preguntas anal\'iticas futuras}, que podr\'ian responderse una vez integrados los datos (no se responden en este avance):
\begin{itemize}
    \item \textquestiondown qu\'e asociaci\'on existe entre la precipitaci\'on acumulada mensual y la presencia o frecuencia de detecciones VIIRS por celda espacial?
    \item \textquestiondown qu\'e celdas espaciales presentan mayor n\'umero o intensidad (FRP) de detecciones en el periodo analizado?
    \item \textquestiondown en qu\'e meses del a\~no se concentra la mayor cantidad de detecciones en el cant\'on Loja?
\end{itemize}

% ==========================================================
% 2. INVENTARIO DE FUENTES
% ==========================================================
\section{Inventario de fuentes}

Se identificaron tres fuentes principales, que combinan datos geoespaciales, observacionales y clim\'aticos.

\begin{center}
\tiny
\setlength{\tabcolsep}{2.2pt}
\renewcommand{\arraystretch}{1.15}
\begin{tabular}{@{}L{1.7cm}L{1.9cm}L{1.5cm}L{1.5cm}L{1.4cm}L{2.7cm}L{2.0cm}L{2.6cm}@{}}
\toprule
\textbf{Fuente} & \textbf{Origen} & \textbf{Tipo} & \textbf{Captura} & \textbf{Formato} & \textbf{Metadatos} & \textbf{Almacena\-miento} & \textbf{Riesgos} \\
\midrule
L\'imite y malla espacial &
INEC, Marco Geoestad\'istico Nacional &
Geoespacial estructurado &
\'Unica vez; disoluci\'on de zonas por c\'odigo DPA &
GeoPackage / GeoJSON &
Capa \texttt{zon\_a}; CRS EPSG:31992$\to$32717; malla 500$\times$500\,m &
PostgreSQL / PostGIS (\texttt{dim\_celda}) &
Geometr\'ia inv\'alida o CRS incorrecto; se valida antes de cargar \\[3pt]
VIIRS (detecciones de incendio) &
NASA FIRMS &
Observacional espacio-temporal &
Autom\'atica (API o descarga por \'area y periodo) &
CSV / JSON &
Sensor, fecha/hora, FRP (MW), indicador de calidad &
MongoDB (detalle) y PostgreSQL (agregado diario) &
Duplicados, falsos positivos, cambios de API/cuota; producto y campo de calidad definitivos a\'un por confirmar \\[3pt]
CHIRPS (precipitaci\'on) &
Climate Hazards Center, UC Santa B\'arbara &
Clim\'atico espacio-temporal &
Autom\'atica (descarga directa) &
CSV / GeoTIFF / Parquet &
Producto CHIRPS v2.0; resoluci\'on $\sim$0,05\textdegree; unidad mm &
PostgreSQL (\texttt{fact\_clima}) y Parquet (respaldo) &
Valores faltantes; cambios en el repositorio de origen \\
\bottomrule
\end{tabular}
\end{center}

\vspace{0.15cm}
\footnotesize
La malla espacial se reconstruye disolviendo las zonas censales del INEC por su c\'odigo de divisi\'on pol\'itico-administrativa (DPA) para obtener el l\'imite cantonal, sobre el cual se ancla una malla sistem\'atica de 500$\times$500\,m. El CRS de trabajo del proyecto es \texttt{EPSG:32717} (WGS84 / UTM 17S); las coordenadas se reproyectan a ese sistema antes de construir cada celda. La arquitectura no utiliza Google Earth Engine. Para VIIRS, el producto/colecci\'on definitivo y el campo de calidad (\texttt{fire\_mask} frente a \texttt{confidence}) todav\'ia deben confirmarse con datos reales. Para CHIRPS, asociar su p\'ixel nativo ($\sim$5\,km) a la malla de 500\,m no aumenta su resoluci\'on real: cada celda hereda por asignaci\'on espacial el valor del p\'ixel que la contiene.
\normalsize

\subsection*{Celdas controladas utilizadas en el prototipo}
\footnotesize
Para demostrar la carga y la consulta del modelo se utilizan dos celdas controladas, verificadas espacialmente contra la capa oficial del INEC:

\begin{center}
\small
\begin{tabular}{@{}lccc c@{}}
\toprule
\textbf{Celda} & \textbf{Longitud} & \textbf{Latitud} & \textbf{Parroquia} & \textbf{Distancia al l\'imite} \\
\midrule
LJ\_TEST\_001 & $-$79,520851 & $-$3,805328 & El Cisne & 1085,00 m \\
LJ\_TEST\_002 & $-$79,516342 & $-$3,809842 & El Cisne & 1066,01 m \\
\bottomrule
\end{tabular}
\end{center}

\vspace{0.1cm}
Las dos celdas est\'an completamente contenidas dentro del cant\'on Loja y de la parroquia El Cisne. Fueron seleccionadas determin\'isticamente para validar el prototipo y no constituyen una muestra representativa del territorio cantonal.
\normalsize

% ==========================================================
% 3. DISEÑO COMPARATIVO SQL vs NoSQL
% ==========================================================
\section{Dise\~no comparativo SQL vs. NoSQL}

\subsection{Entidades}
\textbf{Lado SQL (PostgreSQL/PostGIS):} \texttt{dim\_celda}, \texttt{dim\_fecha}, \texttt{fact\_incendio}, \texttt{fact\_clima}.\\
\textbf{Lado NoSQL (MongoDB):} colecci\'on \texttt{detecciones\_viirs}, con el detalle individual de cada detecci\'on (fecha y hora exactas).

\subsection{Relaciones}

\begin{center}
\footnotesize
\begin{tabular}{@{}lll@{}}
\toprule
\textbf{Tabla con FK} & \textbf{Columnas} & \textbf{Referencia} \\
\midrule
\texttt{fact\_incendio} & \texttt{celda\_id}, \texttt{fecha\_id} & \texttt{dim\_celda}, \texttt{dim\_fecha} \\
\texttt{fact\_clima} & \texttt{celda\_id}, \texttt{fecha\_id} & \texttt{dim\_celda}, \texttt{dim\_fecha} \\
\bottomrule
\end{tabular}
\end{center}

\begin{itemize}
    \item Ambas tablas de hechos tienen \texttt{UNIQUE(celda\_id, fecha\_id)}: una fila por celda y d\'ia.
    \item \texttt{fact\_incendio} y \texttt{fact\_clima} no tienen FK entre s\'i: se cruzan mediante \texttt{JOIN} por \texttt{celda\_id} y \texttt{fecha\_id}.
    \item La relaci\'on entre MongoDB y PostgreSQL es \textbf{l\'ogica}, no una clave for\'anea f\'isica: se resuelve mediante un proceso ETL/ELT que agrupa las detecciones de \texttt{detecciones\_viirs} por \texttt{celda\_id} y \texttt{fecha} para poblar \texttt{fact\_incendio} a grano diario; no existe FK entre ambas bases.
\end{itemize}

\subsection{Datos estables y relacionales (SQL) vs. datos flexibles o anidados (NoSQL)}
\begin{tabular}{@{}lL{13.6cm}@{}}
\toprule
\textbf{SQL} & Malla espacial (geometr\'ia fija), calendario, y registros diarios de incendio/clima por celda: estructura homog\'enea, con integridad referencial (PK/FK) y consultas espaciales de PostGIS. \\
\midrule
\textbf{NoSQL} & Detecciones individuales VIIRS: documentos con coordenadas y calidad anidadas, cuyos campos pueden variar seg\'un la versi\'on del producto satelital sin romper un esquema fijo. \\
\bottomrule
\end{tabular}

\subsection{Justificaci\'on del uso de cada tecnolog\'ia}
\textbf{\textquestiondown Por qu\'e PostgreSQL/PostGIS?} Porque la malla espacial y los hechos diarios de incendio/clima tienen estructura estable y relaciones claras entre celda, fecha y hecho, lo que se beneficia de PK/FK, restricciones y de las capacidades geoespaciales de PostGIS. PostgreSQL resuelve las consultas de agregaci\'on por celda-mes (\texttt{JOIN} + \texttt{GROUP BY}).

\textbf{\textquestiondown Por qu\'e MongoDB?} Porque las detecciones satelitales individuales llegan en JSON semiestructurado y anidado, y su forma puede variar entre versiones del producto sin obligar a migrar un esquema relacional. MongoDB resuelve consultas de agregaci\'on documental sobre el detalle crudo de cada detecci\'on (conteo, FRP promedio o m\'aximo por celda).

\subsection{Consultas que se espera resolver}
\footnotesize
\textbf{En SQL:} n\'umero de detecciones, FRP total y m\'axima por celda-mes; precipitaci\'on acumulada, media y m\'axima por celda-mes, conservando los meses con clima aunque no haya incendio. El procedimiento agrega primero incendio y clima por separado en dos CTE (\texttt{incendio\_mes}, \texttt{clima\_mes}) y luego los combina con un \texttt{LEFT JOIN} desde clima hacia incendio por \texttt{celda\_id + anio + mes}:

\begin{scriptsize}
\begin{verbatim}
WITH incendio_mes AS (            -- agrega detecciones por celda-mes
  SELECT anio, mes, celda_id, SUM(numero_detecciones) AS detecciones,
         SUM(frp_suma_mw) AS frp_total_mw, MAX(frp_maxima_mw) AS frp_maxima_mw
  FROM fact_incendio JOIN dim_fecha USING (fecha_id) GROUP BY anio, mes, celda_id
),
clima_mes AS (                    -- agrega precipitacion por celda-mes
  SELECT anio, mes, celda_id, SUM(precipitacion_diaria_mm) AS precip_acumulada_mm
  FROM fact_clima JOIN dim_fecha USING (fecha_id) GROUP BY anio, mes, celda_id
)
SELECT c.anio, c.mes, c.celda_id,
       COALESCE(i.detecciones, 0) AS detecciones,   -- valido solo si se
       COALESCE(i.frp_total_mw, 0) AS frp_total_mw, -- verifico cobertura VIIRS
       c.precip_acumulada_mm
FROM clima_mes AS c LEFT JOIN incendio_mes AS i USING (anio, mes, celda_id);
\end{verbatim}
\end{scriptsize}

\textbf{En MongoDB:} conteo y FRP promedio/m\'aximo de detecciones por \texttt{celda\_id}, filtrando por el indicador de calidad validada de cada documento (\texttt{calidad.valida}), mediante una canalizaci\'on de agregaci\'on con \texttt{\$match} y \texttt{\$group}.
\normalsize

% ==========================================================
% 4. MODELO PRELIMINAR DE DATOS
% ==========================================================
\section{Modelo preliminar de datos}

\subsection{Tablas principales, claves y granularidad}
\small
\begin{center}
\begin{tabular}{@{}L{2.6cm}L{6.4cm}C{2.1cm}C{3.0cm}@{}}
\toprule
\textbf{Tabla} & \textbf{Descripci\'on} & \textbf{PK} & \textbf{FK} \\
\midrule
\texttt{dim\_celda} & Malla espacial de referencia (una fila por celda) & \texttt{celda\_id} & --- \\
\texttt{dim\_fecha} & Calendario diario del periodo analizado & \texttt{fecha\_id} & --- \\
\texttt{fact\_incendio} & Registro diario de incendio por celda & \texttt{incendio\_id} & \texttt{celda\_id}, \texttt{fecha\_id} \\
\texttt{fact\_clima} & Registro diario de precipitaci\'on por celda & \texttt{clima\_id} & \texttt{celda\_id}, \texttt{fecha\_id} \\
\bottomrule
\end{tabular}
\end{center}
\normalsize

\noindent\footnotesize\textbf{Granularidad:} las tablas de hechos son siempre diarias (\texttt{UNIQUE(celda\_id, fecha\_id)}); la unidad de an\'alisis final, celda-mes, se obtiene agregando esos registros diarios por \texttt{celda\_id + anio + mes} mediante consultas SQL; no existe como filas pre-agregadas en las tablas base.\normalsize

\subsection{Esquema relacional}
\begin{figure}[H]
\centering
\begin{tikzpicture}[
  ent/.style={draw=myNewColorA, thick, rounded corners, fill=myNewColorC!30,
              align=center, minimum width=2.5cm, minimum height=1.05cm, font=\tiny, inner sep=2pt},
  rel/.style={-{Stealth[length=2mm]}, thick, myNewColorB},
  card/.style={font=\scriptsize, myNewColorA}
]
\node[ent] (celda) {\textbf{dim\_celda}\\ PK: \texttt{celda\_id}};
\node[ent, right=1.3cm of celda] (incendio) {\textbf{fact\_incendio}\\ PK: \texttt{incendio\_id}\\ FK: \texttt{celda\_id}, \texttt{fecha\_id}};
\node[ent, right=1.3cm of incendio] (clima) {\textbf{fact\_clima}\\ PK: \texttt{clima\_id}\\ FK: \texttt{celda\_id}, \texttt{fecha\_id}};
\node[ent, right=1.3cm of clima] (fecha) {\textbf{dim\_fecha}\\ PK: \texttt{fecha\_id}};

\draw[rel] (celda) -- (incendio) node[midway, above, card] {1:N};
\draw[rel] (fecha) -- (clima) node[midway, above, card] {1:N};
\draw[rel] (celda.north) to[bend left=30] node[midway, above, card] {1:N} (clima.north);
\draw[rel] (fecha.south) to[bend left=30] node[midway, below, card] {1:N} (incendio.south);
\end{tikzpicture}
\caption{Modelo relacional preliminar: \texttt{dim\_celda} y \texttt{dim\_fecha} se relacionan con \texttt{fact\_incendio} y \texttt{fact\_clima} mediante claves for\'aneas; estas dos \'ultimas se cruzan por \texttt{celda\_id + fecha\_id} sin FK directa entre ellas.}
\end{figure}

\subsection{Ejemplo de documento JSON (NoSQL)}
\footnotesize
La colecci\'on \texttt{detecciones\_viirs} conserva el detalle individual de cada detecci\'on, con fecha y hora exactas; estas detecciones se agregan por \texttt{celda\_id + fecha} para poblar \texttt{fact\_incendio} a grano diario.
\normalsize

\begin{footnotesize}
\begin{verbatim}
{
  "deteccion_id": "V_2023_TEST_001",
  "celda_id": "LJ_TEST_001",
  "fecha": "2023-08-15",
  "hora_utc": "18:42",
  "fuente": "VIIRS",
  "sensor": "VIIRS_SNPP_375m",
  "frp_mw": 18.4,
  "fire_mask": 7,
  "coordenadas": {
    "longitud": -79.520851,
    "latitud": -3.805328
  },
  "calidad": {
    "valida": true
  }
}
\end{verbatim}
\end{footnotesize}

\subsection{Validaci\'on controlada del prototipo}
\footnotesize
La carga y consulta descritas se verificaron con las dos celdas controladas de la secci\'on~2:

\begin{center}
\begin{tabular}{@{}lccc@{}}
\toprule
\textbf{Celda} & \textbf{Detecciones} & \textbf{FRP total / m\'axima} & \textbf{Precipitaci\'on} \\
\midrule
LJ\_TEST\_001 & 1 & 18,4 / 18,4 MW & 12,5 mm \\
LJ\_TEST\_002 & 1 & 9,7 / 9,7 MW & 8,7 mm \\
\bottomrule
\end{tabular}
\end{center}

\vspace{0.1cm}
Se trata de datos controlados: comprueban el funcionamiento del flujo de carga e integraci\'on, no son resultados ambientales representativos y no permiten inferencia alguna para el cant\'on.
\normalsize

% ==========================================================
% 5. ARQUITECTURA PRELIMINAR
% ==========================================================
\section{Arquitectura preliminar}

La figura~\ref{fig:arquitectura} muestra el flujo de datos desde las fuentes hasta el producto anal\'itico.

\begin{figure}[H]
\centering
\begin{tikzpicture}[
  box/.style={draw=myNewColorA, rounded corners, fill=myNewColorC!35,
              align=center, text width=8.2cm, minimum height=0.75cm, font=\small, inner sep=4pt},
  arr/.style={-{Stealth[length=2mm]}, thick, myNewColorA}
]
\node[box] (a) {\textbf{Fuentes} -- malla espacial (INEC), VIIRS (NASA FIRMS), CHIRPS};
\node[box, below=0.22cm of a] (b) {\textbf{Captura e ingesta} -- archivos / API (CSV, JSON, GeoTIFF, Parquet)};
\node[box, below=0.22cm of b] (c) {\textbf{Validaci\'on} -- geometr\'ia, duplicados, valores faltantes};
\node[box, below=0.22cm of c] (e) {\textbf{Almacenamiento} -- MongoDB y PostgreSQL/PostGIS};
\node[box, below=0.22cm of e] (g) {\textbf{Integraci\'on y agregaci\'on} -- clave \texttt{celda\_id + anio + mes}};
\node[box, below=0.22cm of g] (h) {\textbf{Producto anal\'itico} -- dataset celda-mes};
\node[box, below=0.22cm of h] (i) {Estad\'istica / visualizaci\'on / ML};
\draw[arr] (a) -- (b);
\draw[arr] (b) -- (c);
\draw[arr] (c) -- (e);
\draw[arr] (e) -- (g);
\draw[arr] (g) -- (h);
\draw[arr] (h) -- (i);
\end{tikzpicture}
\caption{Arquitectura preliminar de la soluci\'on FireForest.}
\label{fig:arquitectura}
\end{figure}

\noindent\footnotesize La orquestaci\'on de la ingesta, validaci\'on, carga e integraci\'on se prev\'e mediante Airflow; actualmente se encuentra en fase de validaci\'on, no operativa por completo.\normalsize

% ==========================================================
% 6. PREGUNTA ANALÍTICA, UNIDAD DE ANÁLISIS Y DATASET ESPERADO
% ==========================================================
\section{Pregunta anal\'itica, unidad de an\'alisis y dataset esperado}

\textbf{Pregunta anal\'itica:} \textquestiondown qu\'e asociaci\'on existe entre la precipitaci\'on acumulada mensual y la presencia o frecuencia de detecciones VIIRS por celda espacial? Esta pregunta no se responde en este avance; requiere el dataset celda-mes consolidado (secci\'on~7).

\textbf{Unidad de an\'alisis:} celda espacial de 500$\times$500\,m observada durante un mes, con clave \texttt{celda\_id + anio + mes}. Las tablas de hechos permanecen a grano diario; la unidad celda-mes se construye por agregaci\'on.

\textbf{Dataset final esperado:} una tabla con una fila por celda-mes, que combine identificador de celda y coordenadas, n\'umero de detecciones e intensidad de incendio (FRP), y precipitaci\'on acumulada/media/m\'axima del mes.

\vspace{0.15cm}
\noindent\footnotesize\textbf{Cero frente a dato faltante:} un valor \texttt{0} en detecciones o FRP no equivale autom\'aticamente a cobertura VIIRS comprobada; solo puede leerse como ausencia real de detecciones si se verific\'o de forma independiente que la cobertura e ingesta VIIRS de ese mes fueron correctas. \texttt{NULL} representa informaci\'on faltante o no procesada. La prueba de \texttt{LEFT JOIN} con \texttt{COALESCE} descrita en la secci\'on~3.5 valida el comportamiento de la consulta, no constituye evidencia de cobertura satelital real.\normalsize

% ==========================================================
% 7. PLAN PRELIMINAR DE IMPLEMENTACIÓN
% ==========================================================
\section{Plan preliminar de implementaci\'on}

\begin{center}
\scriptsize
\begin{tabular}{@{}C{0.6cm}L{5.0cm}L{5.3cm}C{2.3cm}@{}}
\toprule
\textbf{\#} & \textbf{Actividad} & \textbf{Resultado esperado} & \textbf{Estado} \\
\midrule
1 & Identificaci\'on de fuentes y validaci\'on territorial de las dos celdas controladas & Fuentes documentadas; celdas verificadas dentro del cant\'on Loja & Completada \\
2 & Dise\~no relacional (SQL) & Esquema con PK/FK & Completado \\
3 & Dise\~no NoSQL & Colecci\'on y documento de ejemplo & Completado \\
4 & Carga controlada & Datos cargados en PostgreSQL/PostGIS y MongoDB & Completada \\
5 & Migraci\'on territorial de las celdas controladas & Identificadores y coordenadas actualizados en ambas bases & Completada \\
6 & Adquisici\'on de series reales (VIIRS, CHIRPS) & Datos reales incorporados & Pendiente \\
7 & Malla completa de producci\'on & Malla de 500\,m para todo el cant\'on & Pendiente \\
8 & Cobertura VIIRS & Bandera de cobertura verificada por celda-mes & Pendiente \\
9 & Panel completo celda-mes & Dataset consolidado & Pendiente \\
10 & Automatizaci\'on & Orquestaci\'on de ingesta e integraci\'on (Airflow) & En validaci\'on \\
11 & An\'alisis y producto final & Indicadores, modelo o reporte & Pendiente \\
\bottomrule
\end{tabular}
\end{center}

% ==========================================================
% 8. ALCANCE Y LIMITACIONES
% ==========================================================
\section{Alcance y limitaciones}
\label{sec:alcance}

\begin{itemize}
    \item \textbf{Dominio previsto:} cant\'on Loja.
    \item \textbf{Cobertura actual:} dos celdas controladas verificadas dentro de la parroquia El Cisne, utilizadas exclusivamente para validaci\'on t\'ecnica del prototipo (almacenamiento, reproyecci\'on, agregaci\'on temporal, consultas de integraci\'on SQL/NoSQL).
    \item Este avance corresponde a un \textbf{prototipo acad\'emico}, no a un sistema operativo de alerta temprana de incendios.
    \item Dos celdas dentro de una sola parroquia \textbf{no constituyen una muestra representativa} de la variabilidad ambiental, altitudinal o clim\'atica del cant\'on; los resultados presentados no describen el comportamiento de los incendios del cant\'on Loja.
    \item Las consultas mostradas son ilustrativas: no permiten a\'un an\'alisis inferencial ni predictivo, ni establecer relaci\'on causal entre precipitaci\'on e incendios.
    \item Queda pendiente incorporar series reales completas de VIIRS y CHIRPS, la malla completa de producci\'on y la bandera de cobertura VIIRS.
\end{itemize}

% ==========================================================
% 9. REFERENCIAS
% ==========================================================
\section{Referencias}

\footnotesize
\begin{itemize}
    \item Balaguer-Beser, A., Gonz\'alez, F., Chuncho, C. G., Carri\'on-Paladines, V., \& Juela-Sivisaca, O. (2026). Altitudinal variation in fire behavior in Andean ecosystems in southern Ecuador. \textit{Fire Ecology}, 22, Article 46. \url{https://doi.org/10.1186/s42408-026-00470-y}
    \item Instituto Nacional de Estad\'istica y Censos (INEC). \textit{Marco Geoestad\'istico Nacional}. \url{https://www.ecuadorencifras.gob.ec/documentos/web-inec/Geografia_Estadistica/Micrositio_geoportal/descargas.html}
    \item NASA EOSDIS FIRMS. \textit{VIIRS 375 m Active Fire and Thermal Anomalies (VNP14IMG/VJ114IMG)}. \url{https://firms.modaps.eosdis.nasa.gov}
    \item Funk, C., Peterson, P., Landsfeld, M., Pedreros, D., Verdin, J., Shukla, S., Husak, G., Rowland, J., Harrison, L., Hoell, A., \& Michaelsen, J. (2015). The climate hazards infrared precipitation with stations -- a new environmental record for monitoring extremes. \textit{Scientific Data}, 2, 150066. \url{https://www.chc.ucsb.edu/data/chirps}
    \item PostgreSQL Global Development Group. \textit{PostgreSQL Documentation}. \url{https://www.postgresql.org/docs/}
    \item PostGIS Project Steering Committee. \textit{PostGIS Documentation}. \url{https://postgis.net/documentation/}
    \item MongoDB, Inc. \textit{MongoDB Manual}. \url{https://www.mongodb.com/docs/}
\end{itemize}
\normalsize

\end{document}
